% Capítulo 3 (Principais Trabalhos Relacionados) — Slide 2 — Três lacunas convergentes
\begin{frame}{Três lacunas convergentes}
  \begin{block}{Empírica}
    \small Poucos estudos avaliam políticas de reposição em \highlight{dados
    reais de varejo brasileiro} com demanda altamente intermitente.
  \end{block}
  \begin{block}{Metodológica}
    \small Trabalhos existentes otimizam ou aprendem políticas
    \textit{específicas}, mas não formulam a escolha como
    \highlight{seleção contextual por perfil operacional}.
  \end{block}
  \begin{block}{Operacional}
    \small Pouco explorada a integração, em um \highlight{pipeline
    reprodutível}, entre caracterização, simulação, validação estatística e
    treinamento de um seletor supervisionado.
  \end{block}
  \note{
    A revisão da literatura aponta três lacunas que convergem para a mesma
    direção. A primeira é empírica: faltam estudos com dados reais de varejo
    brasileiro sob demanda altamente intermitente. A segunda é metodológica:
    os trabalhos existentes tendem a otimizar ou treinar uma política
    específica -- calibrar GA, treinar um DQN -- mas não formulam a escolha
    da política como um problema de seleção contextual por perfil
    operacional. A terceira é operacional: ainda é pouco explorada a
    integração, em um único pipeline reprodutível, entre caracterização das
    séries, simulação comparável, validação estatística e treinamento de um
    seletor supervisionado. Essas três lacunas, juntas, motivam diretamente
    a formulação do AIPE.
  }
\end{frame}
